% Ten new terms for five bishop-graph sequences
% LaTeX source; compiles with pdflatex, xelatex, or tectonic.
\documentclass[11pt]{article}

\usepackage[margin=1in]{geometry}
\usepackage{amsmath,amssymb,amsthm}
\usepackage{booktabs}
\usepackage{microtype}
\usepackage[hidelinks]{hyperref}

\newtheorem{theorem}{Theorem}
\newtheorem{lemma}[theorem]{Lemma}
\theoremstyle{remark}
\newtheorem{remark}{Remark}

\newcommand{\seq}[1]{\texttt{#1}}

\title{Ten new terms for five bishop-graph sequences,\\
       via a bishop-to-rook coordinate reduction}
\author{Leo Y. Zhang}
\date{August 2026}

\begin{document}

\maketitle

\begin{abstract}
The $n \times n$ bishop graph is a rook graph in rotated coordinates. We give the
reduction explicitly, prove it, and use it to extend five sequences in the
On-Line Encyclopedia of Integer Sequences by two terms each: \seq{A290719},
\seq{A290769} and \seq{A291595} (connected induced subgraphs, on the black
cells, the white cells, and the whole board) from $n=9$ to $n=11$, and
\seq{A289145} and \seq{A289169} (connected dominating sets) from $n=8$ to
$n=10$. None of the five entries carried a program, a formula or a comment, so
this also supplies a first published method for them.

The reduction turns a count over $O(2^{|B|})$ cell subsets into a sweep over
$O(n)$ classes per side. We separate two regimes sharply: predicates depending
only on \emph{which} classes a set occupies collapse, and reach $n=21$;
predicates needing class \emph{populations} do not, and total domination is
recorded as a negative result. We also show that proper colourings of a bishop
graph are proper edge colourings of a bipartite graph, which by K\"onig's theorem
gives the chromatic number at once but also explains why the associated counting
sequences are, and will remain, short.

Nothing here is a proof of the counts; it is verified computation. Every one of
the ten new terms carries independent confirmation --- seven from a second
algorithm sharing no code path beyond the class grid, three from structural
identities checked against every published term of the sequence they govern.
\end{abstract}

\section{Results}

Write $B$ for the cells of one colour of the $n \times n$ board. The bishop graph
on $B$ joins two cells when they share a diagonal. Table~\ref{tab:results} lists
what was already published and what is added here.

\begin{table}[h]
\centering
\begin{tabular}{llccc}
\toprule
Sequence & Counts & Published to & Now to & New \\
\midrule
\seq{A290719} & connected induced subgraphs, black & $n=9$ & $n=11$ & $+2$ \\
\seq{A290769} & connected induced subgraphs, white & $n=9$ & $n=11$ & $+2$ \\
\seq{A291595} & connected induced subgraphs, bishop & $n=9$ & $n=11$ & $+2$ \\
\seq{A289145} & connected dominating sets, black    & $n=8$ & $n=10$ & $+2$ \\
\seq{A289169} & connected dominating sets, white    & $n=8$ & $n=10$ & $+2$ \\
\bottomrule
\end{tabular}
\caption{Five sequences, ten new terms. ``Published to'' is measured against
each entry's b-file, not its \textsc{data} line; Section~\ref{sec:correction}
explains why that distinction is not pedantry.}
\label{tab:results}
\end{table}

The new values are
\begin{align*}
\seq{A290719}: \quad a(10) &= 1090550900687379, &  a(11) &= 2265142469367980614,\\
\seq{A290769}: \quad a(10) &= 1090550900687379, &  a(11) &= 1134335726831043925,\\
\seq{A291595}: \quad a(10) &= 2181101801374758, &  a(11) &= 3399478196199024539,\\
\seq{A289145}: \quad a(9)  &= 2014079802496,    &  a(10) &= 1073633875253120,\\
\seq{A289169}: \quad a(9)  &= 1011850062768,    &  a(10) &= 1073633875253120.
\end{align*}

Two coincidences in that list are forced rather than accidental, and serve as
checks: at even $n$ a reflection carries the black cells onto the white ones, so
$\seq{A290719}\,a(10) = \seq{A290769}\,a(10)$ and
$\seq{A289145}\,a(10) = \seq{A289169}\,a(10)$; and the bishop graph has exactly
two components, so \seq{A291595} is the sum of the other two throughout.

\section{The reduction}

\begin{theorem}\label{thm:reduction}
Let $B$ be the cells of one colour on the $n \times n$ board and define
\[
  \varphi(r,c) \;=\; \Bigl( \bigl\lfloor (r+c)/2 \bigr\rfloor,\;
                            \bigl\lfloor (r-c+n)/2 \bigr\rfloor \Bigr).
\]
Then $\varphi$ is injective on $B$, and two distinct cells of $B$ share a
diagonal if and only if their images agree in the first coordinate or in the
second.
\end{theorem}

\begin{proof}
Within one colour $r+c$ has fixed parity, hence so does $r-c$. Restricted to
values of a single parity, both $t \mapsto \lfloor t/2 \rfloor$ and
$t \mapsto \lfloor (t+n)/2 \rfloor$ are strictly monotone and therefore
injective. Two cells share an anti-diagonal exactly when $r+c$ agrees, that is
when the first coordinates agree; they share a main diagonal exactly when $r-c$
agrees, that is when the second coordinates agree. Since the pair $(r+c,\,r-c)$
determines $(r,c)$, the map $\varphi$ is injective.
\end{proof}

So the single-colour bishop graph \emph{is} a rook graph on the image cells:
two cells are adjacent exactly when they share an $x$-class or a $y$-class. A
bishop move preserves the parity of $r+c$, so the full bishop graph is the
disjoint union of the black and white graphs and has exactly two components.

The consequence that matters for counting is this. All cells sharing a class are
mutually adjacent, so a set of cells is connected in a rook graph if and only if
the bipartite graph it induces between $x$-classes and $y$-classes is connected.
Since $\varphi$ is injective, each $(x\text{-class}, y\text{-class})$ pair holds
at most one cell, so the class grid is a $0/1$ matrix and a set of cells is
simply a set of occupied grid positions. Counting therefore moves from
$O(2^{|B|})$ subsets of cells to work over $O(n)$ classes per side.

\section{Two algorithms}

\paragraph{Exact-support peeling.}
Let $B(X,Y)$ count the cell subsets whose $x$-support is exactly $X$ and whose
$y$-support is exactly $Y$, obtained from $2^{\mathrm{cells}(X',Y')}$ by
M\"obius inversion over sub-masks, and let $C(X,Y)$ count those that are also
connected. Peeling off the component containing the lowest $x$-class of $X$
gives
\[
  C(X,Y) \;=\; B(X,Y) \;-\!\!\sum_{\substack{X_1 \ni \mathrm{low}(X),\; Y_1 \neq \emptyset \\ (X_1,Y_1) \neq (X,Y)}}\!\! C(X_1,Y_1)\, B(X \setminus X_1,\, Y \setminus Y_1).
\]
Its cost is about $9^n$. It is kept as an independent reference implementation
rather than for production, and that is precisely its value: it shares no code
path with the sweep below beyond the class grid.

\paragraph{Frontier partition DP.}
Sweep the $x$-classes carrying a partition of the $y$-classes into connected
blocks, with label $0$ meaning untouched. At each $x$-class choose a non-empty
subset $T$ of the $y$-classes it meets, weighted by
$\prod_{j \in T} (2^{\mathrm{grid}[i][j]} - 1)$, and merge every block meeting
$T$.

The pruning is what makes this practical. Each $x$-class meets a
\emph{contiguous} interval of $y$-classes, so a $y$-class lying below every
remaining interval can never be touched again. A block confined entirely to that
finalised region can never merge with anything else, so any state carrying such a
stranded block is dropped immediately. This keeps the live state count far below
the Bell number of the frontier width.

\section{Measured cost}

Table~\ref{tab:cost} gives peak live states and wall time on the black board.

\begin{table}[h]
\centering
\begin{tabular}{lrrrrrrrrrrr}
\toprule
$n$ & 1 & 2 & 3 & 4 & 5 & 6 & 7 & 8 & 9 & 10 & 11 \\
\midrule
states & 2 & 2 & 10 & 11 & 101 & 151 & 1596 & 3136 & 39955 & 86321 & 1343614 \\
sec    & 0 & 0 & 0  & 0  & 0   & 0.01 & 0.09 & 0.43 & 9.73 & 53.5  & 1604.7 \\
\bottomrule
\end{tabular}
\caption{Peak live DP states and wall time, black board.}
\label{tab:cost}
\end{table}

The white board reaches $n=11$ in $691.9$\,s, with peak states
$0$, $4$, $4$, $29$, $38$, $380$, $660$, $7529$, $15923$, $225187$, $494815$.
For scale, the peeling
counter needs $510.7$\,s at $n=10$ where the frontier DP needs $53.5$\,s, and
brute force over all $2^{50}$ subsets is not reachable at all.

\begin{remark}[Ceiling]
States grow by a factor of roughly $2.5$ to $15$ per step, so $n=12$ needs
several gigabytes and $n=13$ is out of reach on the machine used here. This work
adds terms; it does not add an order of magnitude. These counts grow doubly
exponentially in $n$, and no algorithm changes that.
\end{remark}

\section{Verification}

Five levels are re-run from cold by the repository's gate, alongside the unit
test suite.

\begin{description}
\item[L0] The reduction of Theorem~\ref{thm:reduction} is re-derived
exhaustively for $n = 1,\dots,12$ in both colours: every pair of cells is checked
for agreement between diagonal-sharing and class-sharing.
\item[L1] The frontier DP is checked against brute-force enumeration over all
cell subsets for $n = 1,\dots,6$, both colours.
\item[L2] Every published term of \seq{A290719} (nine of them) and
\seq{A290769} (eight) is reproduced exactly, indexed by each sequence's true
offset. The offsets are not uniform: \seq{A290719} begins at $n=1$ and
\seq{A290769} at $n=2$, because a $1 \times 1$ board has no white cells.
\item[L3] The frontier DP and the peeling counter agree for $n = 1,\dots,8$ in
both colours, for connected subgraphs and for connected dominating sets.
\item[L4] Structural identities that follow from the problem rather than from the
code: $\mathrm{black}(n) + \mathrm{white}(n) = \mathrm{bishop}(n)$ against all
nine published terms of \seq{A291595}; $\mathrm{black}(n) = \mathrm{white}(n)$
for even $n$ by reflection; and the two differing for odd $n$.
\end{description}

\subsection{Confidence, stated per term rather than in aggregate}

An aggregate claim that ``every term rests on two independent algorithms'' would
have been false when it was first drafted, and the reason is worth recording: the
peeling cross-check then covered two terms of one sequence, while ten terms were
staged across five. Aggregate confidence claims hide exactly that asymmetry.

Stated per term: seven of the ten are confirmed by a \emph{second algorithm}
sharing no code path with the sweep beyond the class grid, and three by a
\emph{structural identity} --- the even-$n$ reflection isomorphism for
\seq{A290769}\,$a(10)$, and $\mathrm{black}+\mathrm{white}=\mathrm{bishop}$ for
\seq{A291595}\,$a(10)$ and $a(11)$ --- each verified against every published term
of the sequence governing it. Neither kind is a re-run of the same code.

\begin{table}[h]
\centering
\begin{tabular}{llrr}
\toprule
Cross-check & Value agreed & Frontier DP & Peeling \\
\midrule
\seq{A289145}\,$a(9)$, black  & 2014079802496    & 8.0\,s  & 147.0\,s \\
\seq{A289169}\,$a(9)$, white  & 1011850062768    & 2.6\,s  & 10.6\,s \\
\seq{A289145}\,$a(10)$, black & 1073633875253120 & 28.1\,s & 528.7\,s \\
\seq{A289169}\,$a(10)$, white & 1073633875253120 & 67.8\,s & 559.0\,s \\
\seq{A290719}\,$a(11)$, black & 2265142469367980614 & 1527.2\,s & 28482.6\,s \\
\seq{A290769}\,$a(11)$, white & 1134335726831043925 & 670\,s & 2737\,s \\
\bottomrule
\end{tabular}
\caption{Independent confirmations by the peeling counter. The two $n=10$
dominating-set values are equal, which the even-$n$ reflection isomorphism
requires and which both algorithms independently produce.}
\label{tab:crosscheck}
\end{table}

\section{A second collapse: domination}

In a rook graph the closed neighbourhood of a cell is the whole of its $x$-class
together with the whole of its $y$-class. A cell is therefore dominated exactly
when its $x$-class or its $y$-class contains a chosen cell.

\begin{lemma}\label{lem:domination}
A set $S$ of cells dominates if and only if the pair
$(\,x\text{-supp}(S),\, y\text{-supp}(S)\,)$ is a vertex cover of the bipartite
graph whose edges are the cells. In particular domination depends only on which
classes $S$ touches, never on which cells it uses inside them.
\end{lemma}

Both this and the neighbourhood description are tested against brute force before
use. The sweep then needs only one bit per live $y$-class recording whether $S$
has touched it, plus one requirement bit, because an $x$-class left empty forces
every $y$-class it meets to be touched.

\subsection{When the reduction pays, and when it does not}

\begin{table}[h]
\centering
\begin{tabular}{@{}llrl@{}}
\toprule
Predicate & What the state must record & States & Reaches \\
\midrule
dominating sets & occupancy, 2 per class & 1125035 & $n=21$ \\
connected induced subgraphs & a partition of the frontier & 1343614 & $n=11$ \\
total dominating sets & population capped at 2, 6 per class & 914611 & $n=11$ \\
\bottomrule
\end{tabular}
\caption{The rule is sharp.}
\label{tab:regimes}
\end{table}

The rule is sharp: \emph{predicates that depend only on which classes are
occupied collapse; predicates that need class populations do not.} Total
domination uses the open neighbourhood, so a chosen cell does not cover itself
and the sweep must distinguish ``one chosen cell here'' from ``two or more''.
That single extra level of state costs seven terms of reach.

Recording a \emph{deficit} rather than a raw requirement matters. Once a class
saturates at two, every requirement on it is already met, so the three states
(count $2$, requirement $0$, $1$ or $2$) behave identically from then on. Merging
them cuts the per-class combinations from nine to six, and $n=10$ from
$36$\,s to $6.6$\,s.

\begin{remark}[Negative result]
Total domination reaches only $n=11$ against the sixteen published terms of
\seq{A303145}, so this method cannot extend it. It is recorded rather than
omitted: the reduction is not a universal win, and knowing which predicates it
fails on is part of the result.
\end{remark}

\subsection{Composing both collapses}

Connected dominating sets need connectivity \emph{and} domination, and the two
compose more cheaply than one might expect. In the connectivity sweep a
$y$-class carries a non-zero block label exactly when it holds a chosen cell,
which is precisely what domination asks about. Domination therefore contributes
one requirement bit per $y$-class rather than a whole extra dimension, and the
combined counter still reaches $n=10$.

There is no full-bishop analogue of that sequence, and there cannot be: a
connected set lies inside one component, so it can never dominate the other. The
OEIS accordingly has black and white variants only.

The two colour components also combine differently depending on the predicate,
which yields independent checks at no cost: connected subgraphs \emph{add},
because a connected subgraph lies in one component, giving
$\seq{A291595} = \seq{A290719} + \seq{A290769}$; dominating sets
\emph{multiply}, because a set dominates a disjoint union exactly when it
dominates each part, giving $\seq{A295898} = \seq{A289164} \cdot \seq{A289170}$.
Both identities are enforced against every published term.

\section{Why the colouring sequences are out of reach}

A cell is an edge between its $x$-class and its $y$-class, and two cells conflict
exactly when they share an endpoint. So proper vertex colourings of the rook
graph are proper \emph{edge} colourings of the bipartite cell graph, and by
K\"onig's edge colouring theorem the chromatic number is the maximum degree:
\[
  \chi(\text{bishop graph}) \;=\; \text{largest diagonal length}.
\]
This is verified for $n = 1,\dots,4$, and counting the proper colourings that use
exactly $\chi$ colours reproduces the published \seq{A371202}
$(1, 2, 12, 528)$.

The same reformulation settles feasibility, negatively. Counting
$\Delta$-edge-colourings of a bipartite graph is, in the regular case, the
problem of counting Latin squares, where only eleven terms are known to anyone. A
sweep would have to carry the full \emph{set} of colours used at each $y$-class,
making the state $(2^{\Delta})^{n_y}$ rather than a constant per class. That
\seq{A371202}, \seq{A371203} and \seq{A371204} stand at five or six terms is not
an accident, and this reduction does not change it. They are recorded here as
understood, not attempted.

\section{A correction, and the process fix it forced}\label{sec:correction}

The dominating-set sequences \seq{A289164}, \seq{A289170} and \seq{A295898} were
also computed here, to $n=21$, and are \emph{not} contributions: their published
b-files already reach $n=50$. The OEIS truncates the \textsc{data} line near 260
characters, so those entries \emph{display} twelve to fifteen terms while holding
fifty. The bishop shortlist had been checked for exactly this; the domination
sequences were picked up later from a grouped listing and were not.

The check is now mechanical rather than remembered: a tool records how far every
published b-file reaches, and a test refuses to let a staged file count as a
contribution unless it exceeds that. Reading the \textsc{data} line as the term
count is the easiest way to lose a day on this kind of work, and it is recorded
here so that the next person does not.

That computation was not wasted. All twenty overlapping terms agree with the
published b-file exactly, including $a(20)$ at 61 digits and $a(21)$ at 67. An
independent author reached $n=50$ by an undisclosed method and the two agree
throughout, which is considerably stronger evidence for the engine than the
fifteen displayed terms it had originally been checked against.

\section{Honest limits}

\begin{itemize}
\item Nothing here is a proof of the counts. It is verified computation, and the
strongest claim available is that two algorithms and five levels of checks
agree.
\item The method is specific to bishop graphs. Local board graphs --- grid, king,
knight --- have a different neighbourhood structure and need a different engine.
\item The fixed-width strip sequences, where a transfer matrix yields a linear
recurrence and thousands of terms, are already well covered in the OEIS by other
contributors; this work deliberately does not compete there.
\item The census of regimes in Table~\ref{tab:regimes} is exhaustive only over
the predicates actually implemented.
\end{itemize}

\section*{Data availability}

The code, the staged b-files and the gate that reproduces every claim above are
available at
\url{https://github.com/Leo-Y-Zhang/GraphRecords}.

\begin{thebibliography}{9}

\bibitem{oeis}
OEIS Foundation Inc., \emph{The On-Line Encyclopedia of Integer Sequences},
\url{https://oeis.org}.

\bibitem{konig}
D.~K\H{o}nig,
\emph{\"Uber Graphen und ihre Anwendung auf Determinantentheorie und
Mengenlehre},
Mathematische Annalen \textbf{77} (1916), 453--465.

\bibitem{stanley}
R.~P. Stanley,
\emph{Enumerative Combinatorics, Volume 1}, 2nd ed.,
Cambridge University Press, 2011.
(M\"obius inversion over the subset lattice.)

\end{thebibliography}

\end{document}
